% Section 1. Introduction.
% Status: first publication-standard rewrite against docs/findings-freeze.md.
% Admitted here: route-only vehicle rotation, excess use, and integration strata.
% Excluded: retired direct-route incidence, preliminary 1.17 hazard, LP returns,
% cross-venue spillover interpretation, and all final cost/persistence inference.

\section{Introduction}
\label{sec:intro}

A vehicle currency is traded not because it is one side of the desired payment, but because it is the market's preferred way to move between two other currencies. This role is highly concentrated in foreign exchange. Theories of vehicle currencies attribute that concentration to a feedback between use and transaction cost: routing more trades through one currency deepens its markets, and deeper markets make it attractive to route still more trades through it \citep{Krugman1980VehicleCurrencies}. Monetary-search models generate a related coordination force because accepting an intermediary object is valuable when others are expected to accept it later \citep{KiyotakiWright1989MoneyMedium}. Related models show how the same coordination forces can sustain a currency's role in invoicing, funding, and balance sheets \citep{GopinathEtAl2020DCP,GopinathStein2021Making,Mukhin2022InternationalPriceSystem}.

The empirical obstacle is that the vehicle route is usually latent. Foreign-exchange records reveal turnover and prices while leaving the market-wide path of a payment and its feasible alternatives unobserved. A currency's international role can therefore be measured by the number of markets quoting it \citep{FlandreauJobst2009Empirics}, or inferred from price and volume responses \citep{Somogyi2026DollarDominanceFX}; realised intermediary use on individual routes remains hidden. This makes the central transition difficult to study. A dominant currency can lose transaction share without losing its intermediary role, and a new currency can gain endpoint demand without becoming the object others trade through.

This paper measures that role in decentralised exchange. A swap transaction records an ordered path through a graph of automated markets. When a route exchanges asset $i$ for asset $o$ through asset $k$, the intermediary is observed directly. The same record identifies the endpoints, so vehicle use can be separated from ordinary demand for $k$ as a source or destination. The setting also follows one currency system from an architecture that required native-asset intermediation to designs in which arbitrary pairs and competing routes are available. That history makes vehicle choice observable. Protocol launches coincide with other market changes, so their dates provide chronology but fail as treatments.

The primary unit is an exact two-leg route, which gives each realised vehicle choice one intermediary and one vote. Vehicle dominance is the share of these choices carried by an asset or asset class. I also construct a vehicle excess-use ratio: an asset's intermediary share divided by its endpoint share on the same route perimeter. A value above one means that the asset's intermediary share exceeds its endpoint share. This distinction is important in a market where stablecoins can become popular endpoints at the same time that they become useful bridges between other assets.

Stablecoins gain sharply as vehicle currencies, although the change is not steady. On calendar days observed in both endpoint years, their equal-weighted share among native and stable vehicles rises from \StableCountBase{} in 2024 to \StableCountEnd{} in 2026, a \StableCountChange{} change with a calendar-HAC standard error of \StableCountSE{}. On value supported within 20\% of route input, the corresponding share rises from \StableValueBase{} to \StableValueEnd{}, a \StableValueChange{} change with a standard error of \StableValueSE{}. Native intermediation rebounds in 2023 and 2024; stablecoins recover the value lead in 2025, while their lead by route count remains short-lived. I therefore describe the episode as a rotation. Persistent replacement would require a longer post-change lead.
% docs/findings-freeze.md, vehicle-dominance claim; exact two-leg unit

An aggregate share leaves open whether stablecoins replace the native asset on the same endpoint pairs or gain as routed activity moves across pairs. I distinguish these cases with a five-term accounting of pooled route count, whose stable share rises from \MarketBridgeBase{} to \MarketBridgeEnd{}. Pairs represented in the admitted route-component ledger in only one endpoint year contribute \MarketSupportBridge{} of the \MarketBridgeTotal{} increase. Among pairs represented in both endpoint years, differences in whether a native-or-stable vehicle route appears contribute \VehicleRoleSupportBridge{}. For pairs carrying these vehicle routes in both years, changes in their shares of admitted route-component activity contribute \MarketActivityReweight{}, changes in realised vehicle incidence contribute \VehicleIncidenceReweight{}, and changes in stable share contribute \WithinPairStableShare{}. Only the last term measures stable-for-native substitution within pairs carrying the vehicle role in both years. Most of the observed increase accompanies changes in pair representation, activity, and vehicle use.
% output/exhibits/vehicle_transition_pair_decomposition.jsonl; producer e16796c; evidence commit 23ee3a2

The stablecoin category also conceals an important difference between USDC and USDT. Together they account for \JointStableContribution{} of the increase in stable count share, but USDC begins the comparison with disproportionate vehicle use and changes little. USDT's count excess-use ratio rises from \USDTCountExcessBase{} to \USDTCountExcessEnd{}, while its strict-value ratio rises from \USDTValueExcessBase{} to \USDTValueExcessEnd{}. Its strict-value intermediary-minus-endpoint share moves from \USDTEndpointGapBase{} to \USDTEndpointGapEnd{}. USDT thus becomes more prominent in the intermediary position even relative to its use as an endpoint. The present comparisons do not distinguish among its backing, redemption, liquidity, and venue-reach features.
% docs/findings-freeze.md, token decomposition and excess-use transition

A widening trading network could produce these changes without a stronger coordination role for stablecoins. Cross-venue execution expands rapidly over the sample, yet the share of economic routes using a true intermediary falls from \FullIntermediationBase{} in 2022 to \FullIntermediationEnd{} in 2026 on the full perimeter and to \BalancedIntermediationEnd{} on a balanced five-venue perimeter. Among routes that remain indirect, stablecoin share rises in every observed integration-by-complexity count cell and rises \CrossVenueRotationPremium{} more across venues than within one venue, with a standard error of \CrossVenueRotationSE{}. Changes in broad route types therefore cannot account for the entire rotation. These comparisons do not hold fixed the endpoint pair, trade size, feasible paths, or router's decision state, so better search and changing venue reach remain viable explanations.
% output/exhibits/cross_venue_routing_inference.jsonl;
% docs/findings-freeze.md, integration-by-complexity and paired-date interaction

Protocol launches do not supply the missing fixed comparison. Entry and exit windows alter the set of observed routes, and role appearance or disappearance can coincide mechanically with the event window. I use these dates to describe chronology and do not treat them as estimates of an architecture effect.
% output/exhibits/architecture_transition_support.jsonl;
% output/exhibits/architecture_role_margin_support.jsonl

The paper makes three contributions. First, it observes the vehicle on each realised route and therefore separates vehicle status, which is binary on a route, from vehicle dominance, which is a continuous market share. The distinction turns the making of a vehicle currency from a quotation or aggregate-turnover object into a sequence of realised intermediary choices.

Second, the vehicle excess-use ratio benchmarks an asset's intermediary share against its observed endpoint share on the same route perimeter. It therefore distinguishes growth in ordinary stablecoin use from disproportionate use between other assets without requiring an exclusion restriction or a structural model.

Third, the paper distinguishes a change in vehicle choice within an endpoint pair from changes in which pairs carry vehicle routes, how much admitted route-component activity they receive, and how often that activity uses an intermediary. This distinction links aggregate vehicle dominance to the cross-section of routed markets instead of treating it as the choice of a representative pair.

The remainder of the paper proceeds as follows. Section~\ref{sec:setting} defines routes, vehicle dominance, and the data perimeter. Section~\ref{sec:dominance} measures the rotation and vehicle excess use. Section~\ref{sec:survival} defines persistence and the data required to estimate turnover. Section~\ref{sec:rivals} evaluates routing integration and the competing mechanisms. Section~\ref{sec:executability} describes exact-state reconstruction and the support required for counterfactual route costs. Section~\ref{sec:conclusion} concludes.
